% 73-09(73쪽) 보기 — 닫힌구간 [-1,1] 에서의 y=f(x) 그래프 네 개(ㄱ 포물선 · ㄴ 꺾은선 · ㄷ 곡선 · ㄹ 한 점에서 끊긴 곡선).
% 스캔 화질이 낮아 TikZ 로 다시 그림 — 라벨은 원문 그대로(-1, O, 1, x, y, y=f(x))만 두고, ㄹ 의 빈 점(○)·채운 점(●)은 원문 그대로 옮겼다.
\begin{tikzpicture}[x=0.5cm, y=0.5cm, line width=0.5pt, font=\footnotesize]
  % ㄱ. 꼭짓점이 원점인 포물선
  \begin{scope}
    \node at (-2.5,1.8) {ㄱ.};
    \draw[->, >=stealth, line width=0.7pt] (-1.55,0) -- (1.95,0) node[below=1pt] {$x$};
    \draw[->, >=stealth, line width=0.7pt] (0,-0.75) -- (0,2.25) node[left=1pt] {$y$};
    \node[below left=-2pt and -3pt] at (0,0) {$\pt{O}$};
    \draw[densely dashed, line width=0.4pt] (-1,1.35) -- (1,1.35);
    \draw[dotted, line width=0.4pt] (-1,1.35) -- (-1,0);
    \draw[dotted, line width=0.4pt] (1,1.35) -- (1,0);
    \draw[line width=0.6pt, domain=-1:1, samples=70, smooth] plot (\x, {1.35*\x*\x});
    \node[below=1pt] at (-1,0) {$-1$};
    \node[below=1pt] at (1,0) {$1$};
    \node[anchor=west] at (0.18,1.95) {$y=f(x)$};
  \end{scope}
  % ㄴ. 꼭짓점에서 꺾인 이등변삼각형 꼴
  \begin{scope}[shift={(7.8,0)}]
    \node at (-2.5,1.8) {ㄴ.};
    \draw[->, >=stealth, line width=0.7pt] (-1.55,0) -- (1.95,0) node[below=1pt] {$x$};
    \draw[->, >=stealth, line width=0.7pt] (0,-0.75) -- (0,2.25) node[left=1pt] {$y$};
    \node[below left=-2pt and -3pt] at (0,0) {$\pt{O}$};
    \draw[line width=0.6pt] (-1,0) -- (0,1.5) -- (1,0);
    \node[below=1pt] at (-1,0) {$-1$};
    \node[below=1pt] at (1,0) {$1$};
    \node[anchor=west] at (0.15,1.3) {$y=f(x)$};
  \end{scope}
  % ㄷ. 극대 뒤 원점을 지나 극소로 내려가는 곡선
  \begin{scope}[shift={(0,-4.7)}]
    \node at (-2.5,1.8) {ㄷ.};
    \draw[->, >=stealth, line width=0.7pt] (-1.55,0) -- (1.95,0) node[below=1pt] {$x$};
    \draw[->, >=stealth, line width=0.7pt] (0,-1.75) -- (0,1.65) node[left=1pt] {$y$};
    \node[below left=-2pt and -3pt] at (0,0) {$\pt{O}$};
    \draw[dotted, line width=0.4pt] (-1,0.6) -- (-1,0);
    \draw[dotted, line width=0.4pt] (1,-0.6) -- (1,0);
    \draw[line width=0.6pt, domain=-1:1, samples=90, smooth] plot (\x, {-1.05*sin(145.2*\x)});
    \node[below=1pt] at (-1,0) {$-1$};
    \node[above=1pt] at (1,0) {$1$};
    \node[anchor=west] at (0.2,-1.5) {$y=f(x)$};
  \end{scope}
  % ㄹ. 한 점에서 값만 따로 떨어진 증가 곡선
  \begin{scope}[shift={(7.8,-4.7)}]
    \node at (-2.5,1.8) {ㄹ.};
    \draw[->, >=stealth, line width=0.7pt] (-1.55,0) -- (1.95,0) node[below=1pt] {$x$};
    \draw[->, >=stealth, line width=0.7pt] (0,-0.75) -- (0,2.25) node[left=1pt] {$y$};
    \node[below left=-2pt and -3pt] at (0,0) {$\pt{O}$};
    \draw[densely dashed, line width=0.4pt] (-1,0.45) -- (-1,0);
    \draw[densely dashed, line width=0.4pt] (1,1.45) -- (1,0);
    \draw[dotted, line width=0.4pt] (-0.5,0.95) -- (-0.5,0);
    \draw[line width=0.6pt, domain=-1:1, samples=90, smooth] plot (\x, {0.45+sqrt((\x+1)/2)});
    \draw[fill=white, line width=0.5pt] (-0.5,0.95) circle (2pt);
    \fill (-0.5,0) circle (2pt);
    \node[below=1pt] at (-1,0) {$-1$};
    \node[below=1pt] at (1,0) {$1$};
    \node[anchor=west] at (0.18,1.95) {$y=f(x)$};
  \end{scope}
\end{tikzpicture}
